\documentclass[12pt, letterpaper]{article}

%----------------------------------------------------------------%
% Preamble: Packages and Document Setup (Match Calculus Reports)
%----------------------------------------------------------------%

\usepackage{amsmath, amssymb}
\usepackage[margin=1in]{geometry}
\usepackage{graphicx}
\usepackage[hidelinks]{hyperref}
\hypersetup{colorlinks=false, pdfborder={0 0 0}}
\usepackage{xcolor}
\usepackage{listings}
\usepackage{float}
\usepackage{booktabs}

% Define the guidance text color
\definecolor{guidancecolor}{rgb}{0.5, 0.5, 0.5}

% Python code styling for the appendix
\lstdefinestyle{pythonstyle}{
    commentstyle=\color{black!60},
    keywordstyle=\color{blue!60!black},
    numberstyle=\tiny\color{black!60},
    stringstyle=\color{green!40!black},
    basicstyle=\ttfamily\footnotesize,
    breaklines=true,
    frame=single,
    rulecolor=\color{black!20},
    backgroundcolor=\color{black!2},
    tabsize=2
}

% Helper: include figure if present, otherwise show a placeholder box.
\newcommand{\includeorplaceholder}[1]{%
\IfFileExists{#1}{%
    \includegraphics[width=0.85\linewidth]{#1}%
}{%
    \fbox{\parbox[c][0.33\textheight][c]{0.85\linewidth}{\centering
    \texttt{\detokenize{#1}} not found.\\
    Place the generated figure next to this \texttt{.tex} file.}}%
}%
}

\title{\vspace*{2cm}\Huge \textbf{Math4AI: Probability \& Statistics}\\[0.3cm]
\LARGE Assignment 7 Report: Markov Chains \& Sampling Methods\\[0.2cm]
\large National AI Academy}
\author{\textbf{[Fərman Xankişiyev]} \\ \texttt{[ferman.xankisiyev25@aiacademy.az]}}
\date{February 26, 2026 - March 17, 2026}

\begin{document}

\maketitle
\newpage
\tableofcontents
\newpage

%================================================================%
\section{Task 7.1: PageRank (Power Iteration)}
%================================================================%

\textit{\textcolor{guidancecolor}{
[Guidance: Explain how PageRank treats the web as a Markov Chain. Briefly describe the update rule and the role of the damping factor $d$.]
}}

\subsection{Results}
\textit{\textcolor{guidancecolor}{
[Guidance: Report the final scores. Which node had the highest score? Explain why a node with fewer incoming links might still rank higher than a node with many incoming links (Hint: "Vote Quality").]
}}

\begin{itemize}
    \item Node 0: \texttt{[0.1281]}
    \item Node 1: \texttt{[0.1191]}
    \item Node 2: \texttt{[0.0974]}
    \item Node 3: \texttt{[0.1233]}
    \item Node 4: \texttt{[0.0778]}
    \item Node 5: \texttt{[0.0985]}
    \item Node 6: \texttt{[0.0838]}
    \item Node 7: \texttt{[0.1195]}
    \item Node 8: \texttt{[0.0965]}
    \item Node 9: \texttt{[0.0560]}
    \item \textbf{Most Important Node: 0} \texttt{[score: 0.1281]}
\end{itemize}

\subsection{Graph Visualization}
\begin{figure}[H]
\centering
\includeorplaceholder{pagerank_graph.png}
\caption{Visualization of the Mini-Internet graph structure.}
\end{figure}

%================================================================%
\section{Task 7.2: Metropolis-Hastings (MCMC)}
%================================================================%

\textit{\textcolor{guidancecolor}{
[Guidance: Explain the concept of MCMC. How does the acceptance ratio $A(x \to x')$ ensures we sample from the target density?]
}}

\subsection{Trace Analysis \& Histogram}
\textit{\textcolor{guidancecolor}{
[Guidance: A good MCMC chain should look like a "fuzzy caterpillar." Did your chain successfully jump between the two modes (-4 and +4), or did it get stuck? Discuss the concept of "mixing." check the \texttt{mh\_trace\_hist.png} output.]
}}

\begin{figure}[H]
\centering
\includeorplaceholder{mh_trace_hist.png}
\caption{Left: Trace plot showing the "mixing" of the chain. Right: Histogram of samples vs. True Density.}
\end{figure}

%================================================================%
\section{Task 7.3: Gibbs Sampling}
%================================================================%

\textit{\textcolor{guidancecolor}{
[Guidance: Describe how Gibbs sampling differs from Metropolis-Hastings (i.e., we accept every sample, but move one coordinate at a time).]
}}

\subsection{Reconstructing the Joint Distribution}
\textit{\textcolor{guidancecolor}{
[Guidance: The scatter plot should show a diagonal ellipse. Explain how we achieved this 2D shape while only ever sampling from simple 1D Normal conditionals. Why is this useful in high dimensions?]
}}

\begin{figure}[H]
\centering
\includeorplaceholder{gibbs_scatter.png}
\caption{2D Samples generated via Gibbs Sampling ($\rho=0.8$).}
\end{figure}

%================================================================%
\section{AI Declaration}
%================================================================%
\textit{\textcolor{guidancecolor}{
[Guidance: If you used AI tools, describe exactly what you used them for (e.g., debugging, explanation, drafting). If not, write ``No AI tools were used.'' ]
}}

%================================================================%
\appendix
\section{Python Implementation}
%================================================================%

\textit{\textcolor{guidancecolor}{
[Guidance: Paste your completed Python code here exactly as in your submitted \texttt{.py} file.]
}}

\begin{lstlisting}[style=pythonstyle, caption={Full Python source code for Assignment 7.}, label={lst:code}]
# --- Student code pasted here ---
\end{lstlisting}

\end{document}
